\documentclass[11pt,a4paper]{article}
\usepackage[T2A]{fontenc}
\usepackage[utf8]{inputenc}
\usepackage[russian,english]{babel}
\usepackage{amsmath,amssymb,amsthm,mathtools}
\usepackage{setspace}
\usepackage{enumitem}
\usepackage[hidelinks]{hyperref}
\setstretch{1.3}
\frenchspacing
\hyphenpenalty=10000
\renewcommand{\leq}{\leqslant}
\renewcommand{\le}{\leqslant}
\renewcommand{\geq}{\geqslant}
\renewcommand{\ge}{\geqslant}
\newcommand{\head}[1]{\par\vspace{15pt}\noindent\textbf{\underline{#1}}\quad\ignorespaces}
\newcommand{\src}[1]{\unskip\nobreak\hfil\penalty50\hskip1em\hbox{}\nobreak\hfill(#1){\parfillskip=0pt\par}}
\begin{document}
\begin{center}
{\huge\bfseries Methods of proof -- 4\par}\vspace{0.8em}
{\Large Invariants, monovariants and double counting\par}\vspace{0.8em}
{\foreignlanguage{russian}{13 сентября 2026 года}\par}
\end{center}\vspace{-3.3pt}

\noindent\textbf{\underline{Theoretical minimum.}}\par\vspace{10pt}
\noindent\textbf{\underline{Fact 1 (invariant).}}\quad A quantity unchanged by the allowed
operations separates the reachable states from the unreachable ones. The classical discrete
invariants are parity, a residue, a sum with weights, the sign of a permutation, a colouring; the
continuous ones are a first integral of a differential equation and a conserved energy
(Functional and differential equations -- 3); the algebraic ones are the rank, the determinant,
the trace and the characteristic polynomial, which do not change under similarity (Linear
algebra -- 3).

\head{Fact 2 (monovariant).} A quantity that changes monotonically proves that a process
terminates or that a configuration stabilises; a bounded integer monovariant gives termination,
and a Lyapunov function is its continuous counterpart. The two are used in the same way: find the
quantity, prove monotonicity, look at what happens at the end.

\head{Theorem 1 (double counting).} Counting the same set in two ways gives an identity; counting
incidences in two ways gives an inequality. The systematic form is $\sum_v\deg v=2|E|$ for graphs
and $\sum_i|A_i|=\sum_x\#\{i:x\in A_i\}$ for set systems, and the linear-algebraic strengthening
is the rank bound of Linear algebra -- 6.

\head{Fact 3 (a generating function is a ledger).} Encoding a family by a polynomial or a power
series turns a counting statement into an algebraic identity; the coefficient extraction by roots
of unity is in Trig-2, and the analytic machinery in Series and infinite products -- 4.

\medskip
\noindent\rule{\textwidth}{0.4pt}

\begin{enumerate}[itemsep=0.4em,topsep=7.4pt]
\item Prove that the 15-puzzle position obtained by swapping two neighbouring tiles is
unreachable, by using the sign of the permutation together with the parity of the position of the
empty cell, and formulate the general principle for a group acting on configurations.

\item On a board the numbers may be replaced by $a,b\mapsto\frac{a+b}{\sqrt2},\frac{a-b}{\sqrt2}$.
Find an invariant, and decide which multisets are reachable from a given one when there are only
two numbers on the board; then, for the operation $a,b\mapsto a+b,\ ab$ on positive reals, find a
monovariant and deduce that no configuration can occur twice.

\item Prove that the process in which two numbers $a\ge b$ are replaced by $a-b$ and $b$
terminates, by exhibiting a monovariant; deduce the Euclidean algorithm and its complexity, and
give a continuous analogue where the monovariant is an integral.

\item Prove $\sum_v\deg v=2|E|$ and use double counting to prove that a graph on $n$ vertices with
no $4$-cycle has at most $\frac12\bigl(1+\sqrt{4n-3}\bigr)\cdot\frac n2$ edges; then prove the
Oddtown theorem by the rank bound and compare the two methods (Linear algebra -- 6).

\item (Putnam 2018) Let $S$ be the set of sequences of length 2018 whose terms are in
$\{1,2,3,4,5,6,10\}$ and sum to 3860. Prove that the cardinality of $S$ is at most
$2^{3860}\cdot\binom{2018}{2048}^{2018}$.
\end{enumerate}
\end{document}
